\section{Mortal Sin and the Life of Grace}

The danger of hell concerns our relation to God now. Mortal sin destroys charity in the heart and deprives the soul of sanctifying grace. It turns the person away from God as his last end. Venial sin wounds charity while leaving that friendship intact. The difference concerns the life of the soul, and an honest examination must preserve it (CCC 1854--1856).

St.~Thomas Aquinas explains the gravity of mortal sin by considering the principle which orders our will to God. That principle is adherence to our last end through charity. When sin destroys it, the sinner cannot repair the supernatural loss by his own power. The damage is irreparable from the resources of the fallen person himself, though God can restore what sin has destroyed (ST I--II, q.~87, a.~3).

The last end is the good for whose sake all other goods must be ordered. A Christian may affirm that God is his end while treating a particular desire as exempt from that ordering. A person who will knowingly commit a grave injustice in order to retain wealth has, in that choice, preferred the lesser good to God's law and friendship. The whole life need not be outwardly disorderly for such a choice to be serious. A single deliberate refusal in grave matter can contradict the charity by which the soul belongs to God.

Nor does the fact that the sinful act has finished mean that the relation has healed itself. A lie can be over as a spoken event while its harm remains. A betrayal can cease as an outward act while the betrayer still refuses repentance. Mortal sin demands a return of the will to God and the forgiveness he alone gives. Time can make the memory less vivid without making the sinner reconciled. We therefore need to examine what has been repented of, as well as what has ceased to trouble us.

Three conditions must be present together for a personal sin to be mortal: grave matter, full knowledge, and deliberate consent (CCC 1857). Grave matter concerns the act's seriousness; knowledge and consent concern the person's responsible choice. This protects the examination from both carelessness and confusion. We must consider the real evil of an act and the freedom with which we chose it. An alarming feeling by itself settles neither question.

Temptation is not consent. A thought can present itself against our will; an attraction can be felt while the will refuses it; fear or agitation can persist after a faithful refusal. The presence of struggle is therefore no proof that grace has been lost. The question is what we knowingly and freely did with the temptation. In the same way, having little emotional disturbance after a grave choice does not turn that choice into innocence. Conscience must be formed by truth, not by the intensity or absence of feeling.

Ignorance, duress, fear, habit, and other impediments can diminish responsibility; deliberate cultivation of ignorance and hardness of heart do not excuse it (CCC 1735, 1859--1860). God knows the whole act and the person who performs it. An examination of conscience should be candid about circumstances without making them a shelter for dishonesty. Where the conscience remains troubled and uncertain, the counsel of a steady confessor helps it attend to what is real. Repeatedly reopening settled doubts can distract from the acts of trust and obedience which are presently required.

Venial sin also deserves repentance. It leaves charity alive but weakens its exercise, nourishes disordered attachment, and can dispose the soul toward more serious consent (CCC 1862--1863). Its danger is often the ease with which it becomes cherished. We begin to regard a particular unkindness, vanity, indulgence, or evasion as something we are entitled to keep. The annual examination can uncover this consent to remain unchanged, even where no grave sin has occurred.

The distinction between mortal and venial sin gives proportion to that examination. It does not make venial sin desirable, nor does it turn every imperfection into eternal catastrophe. A person who loves God wants to be rid of whatever offends him. The motive can mature from avoiding destruction to refusing even a smaller injury to friendship.

An examination becomes more truthful when it follows one defended fault through its excuses and effects. We may have admitted impatience in general while refusing to acknowledge the particular person we repeatedly humiliate. Naming the conduct makes repentance less evasive. It also helps us recognize grace at work when the same provocation no longer draws the same willing cruelty. The annual review can then include thanksgiving for a change we did not produce unaided. Sorrow and gratitude both require attention to what actually happened, rather than a general judgment that we have been either good or bad.

The soul that discovers mortal sin is being called to return. Aquinas's account of the lost principle does not close the door upon the living sinner; it tells him where help must come from. We cannot bargain ourselves back into charity or erase the past by becoming impressive in another part of life. We can receive Christ's mercy, consent to the conversion he gives, and submit the wound to the sacrament he has established.

